% !PRINTED 257
\weberpage{257}

Wenn nun $n$ gr\"osser als 1 ist, so kann man $\varrho$ und $1+\varrho$ von Null verschieden annehmen und dann $\sigma=-1:(1+\varrho)$ setzen, wodurch die Substitution (9) in $A_0$ \"ubergeht.

Ist aber $n=1$ und $p=2$, so tritt der \gesperrt{erste Ausnahmefall} ein; denn dann ist immer eine der beiden Zahlen $\varrho$ und $1+\varrho$ gleich Null. Dieser Fall f\"uhrt auf eine Gruppe $6^{\text{ten}}$ Grades
\[
E = \begin{pmatrix}1,&0\\0,&1\end{pmatrix}, \begin{pmatrix}1,&0\\1,&1\end{pmatrix}, \begin{pmatrix}0,&1\\1,&1\end{pmatrix}, \begin{pmatrix}1,&1\\1,&0\end{pmatrix}, \begin{pmatrix}1,&1\\0,&1\end{pmatrix}, \begin{pmatrix}0,&1\\1,&0\end{pmatrix},
\]
in der ein Normaltheiler $3^{\text{ten}}$ Grades enthalten ist:
\[
G = \begin{pmatrix}1,&0\\0,&1\end{pmatrix}, \begin{pmatrix}0,&1\\1,&1\end{pmatrix}, \begin{pmatrix}1,&1\\1,&0\end{pmatrix}.
\]

Von diesem Ausnahmefalle sehen wir jetzt ab. Wir haben dann weiter:

4) Die Gruppe $G$ enth\"alt jede Substitution von der Form
\weq{10}{\begin{pmatrix} 1, & 0 \\ \xi, & 1 \end{pmatrix},}
worin $\xi$ eine beliebige Zahl in $\mathfrak{E}$ ist, ausgenommen in dem Falle $n=1$, $p=3$.

Denn nach 3) enth\"alt $G$ die Substitution
\[
\begin{pmatrix} 0, & \varrho \\ -\varrho^{-1}, & 0 \end{pmatrix}\begin{pmatrix} 0, & 1 \\ -1, & 0 \end{pmatrix}\begin{pmatrix} 0, & -\varrho \\ \varrho^{-1}, & 0 \end{pmatrix}\begin{pmatrix} 0, & 1 \\ -1, & 0 \end{pmatrix} = \begin{pmatrix} \varrho^2, & 0 \\ 0, & \varrho^{-2} \end{pmatrix},
\]
wenn $\varrho$ eine beliebige von Null verschiedene Zahl in $\mathfrak{E}$ ist. Daraus folgt, dass auch
\weq{11}{\begin{pmatrix} 1, & 0 \\ -\lambda, & 1 \end{pmatrix}\begin{pmatrix} \varrho^{-2}, & 0 \\ 0, & \varrho^2 \end{pmatrix}\begin{pmatrix} 1, & 0 \\ \lambda, & 1 \end{pmatrix}\begin{pmatrix} \varrho^2, & 0 \\ 0, & \varrho^{-2} \end{pmatrix} = \begin{pmatrix} 1, & 0 \\ \lambda(\varrho^4-1), & 1 \end{pmatrix}}
f\"ur jedes beliebige $\lambda$ in $G$ vorkommt. Kann man nun $\varrho$ von Null verschieden so annehmen, dass $\varrho^4-1$ nicht verschwindet, so kann man f\"ur jedes $\xi$ die Zahl $\lambda$ aus $\lambda(\varrho^4-1)=\xi$ bestimmen, und erh\"alt also aus (11) jede Substitution der Form (10). Um die M\"oglichkeit hiervon zu beurtheilen, nehmen wir eine primitive Wurzel $\gamma$ von $\mathfrak{E}$ an und setzen $\varrho=\gamma^h$.

Kann man den Exponenten $h$ so annehmen, dass $4h$ nicht durch $p^n-1$ theilbar ist, so gen\"ugt $\varrho$ unserer Forderung. Dies ist immer m\"oglich, wenn $p=2$, $n>1$ ist, und wenn $p^n-1>4$ ist, und es bleiben also nur die beiden F\"alle $n=1$, $p=3$ und $n=1$, $p=5$ noch zweifelhaft.
% [signature line "Weber, Algebra. II.  17" omitted -- gathering signature, not body text]
